\documentclass[11pt,a4paper]{article}

\usepackage[utf8]{inputenc}
\usepackage[margin=1in]{geometry}
\usepackage{amsmath,amssymb,amsthm}
\usepackage{array}
\usepackage{hyperref}

\hypersetup{
  colorlinks=true,
  linkcolor=blue,
  citecolor=blue,
  urlcolor=blue
}

\newtheorem{theorem}{Theorem}[section]
\newtheorem{proposition}[theorem]{Proposition}
\newtheorem{lemma}[theorem]{Lemma}
\newtheorem{corollary}[theorem]{Corollary}
\theoremstyle{definition}
\newtheorem{definition}[theorem]{Definition}
\theoremstyle{remark}
\newtheorem{remark}[theorem]{Remark}

\newcommand{\qbinom}[3]{\begin{bmatrix} #1 \\ #2 \end{bmatrix}_{#3}}
\newcommand{\lcm}{\operatorname{lcm}}
\newcommand{\carry}{\operatorname{carry}}

\title{\textbf{Disproof of Conjecture 00000007675:\\
Vanishing Valuation of Gaussian Binomials at Prime Powers\\
and Failure of the Conjectured LCM Criterion}}
\author{Formal Mathematical Resolution \& Core Lean 4.19.0 Verification}
\date{November 2026}

\begin{document}
\maketitle

\begin{abstract}
We resolve \textbf{Conjecture 00000007675} in the negative (\textbf{DISPROVED}).
First, we prove from the canonical $q$-Pascal, $q$-factorial, and multiplicative
definitions of the Gaussian binomial coefficient $\qbinom{n}{k}{q}$ that
$\qbinom{n}{k}{q} \equiv 1 \pmod{q}$ for every integer $q \ge 2$ and all
$0 \le k \le n$. Consequently, for every prime $p$, every integer $r \ge 1$, and
every $0 \le k \le n$, the $p$-adic valuation satisfies
$v_p\!\left(\qbinom{n}{k}{p^r}\right) = 0$, and in particular
$v_p\!\left(\qbinom{n}{k}{p}\right) = 0$ identically.
Second, we analyze the first clause of the conjecture by proving equivalence between
schoolbook base-$p$ ripple carries and modulo-$p^j$ residue carries, together with
Kummer's exact termwise identity
$\lfloor (a+b)/p^j \rfloor - \lfloor a/p^j \rfloor - \lfloor b/p^j \rfloor = \carry_j(a,b)$;
under standard $1$-indexed carries this reduces the first clause to a tautological
$0 = 0$ cancellation, whereas under $0$-indexed carries it fails already at
$(p,n,k) = (2,2,1)$.
Third, we decisively refute the second clause---asserting that the alternative
LCM-expression $\lfloor \log_p(\lcm(k+1,\dots,n)/(k!(n-k)!))\rfloor$ equals
$v_p\!\left(\qbinom{n}{k}{p}\right)$ \emph{only when} $p \le n - k$---across all eight
natural syntactic and arithmetic parsings of the LCM-expression and under all logical
readings (\emph{only when}, \emph{when}, and \emph{if and only if}), even when
restricted to strictly interior parameters $0 < k < n$ and even for every individual
prime $p$. All definitions, canonical specifications, infinite counterexample
families, and master disproof theorems are formalized from first principles in Core
Lean~4.19.0 without external libraries or non-standard axioms.
\end{abstract}

\section{Official Statement and Executive Summary}

Conjecture~00000007675 reads verbatim in English as follows:

\begin{quote}
\textbf{Definition:} The valuation $v_p([n\text{ choose }k]_q)$ at a prime power is
defined by substituting $q=p$ and then taking the $p$-adic valuation.
\textbf{Conjecture:}
\[
  v_p\!\left(\qbinom{n}{k}{p}\right)
  = \sum_{j \ge 1} \left(
      \left\lfloor \frac{n}{p^j} \right\rfloor
      - \left\lfloor \frac{k}{p^j} \right\rfloor
      - \left\lfloor \frac{n-k}{p^j} \right\rfloor
      - \carry_j(k,n-k)
    \right),
\]
where $\carry_j$ counts the $j$-th carry; and the alternative LCM-expression holds
only when $p \le n - k$. (\emph{carry formula for Gaussian binomial valuations})
\end{quote}

\noindent
The accompanying official Chinese statement supplies the explicit form of the
``alternative LCM-expression'', stating that the valuation equals a
\[
  \left\lfloor \log_p \frac{\lcm(k+1,\dots,n)}{k!(n-k)!} \right\rfloor
\]
type expression \emph{only when} (\emph{jin dang}) $p \le n - k$.

\paragraph{Verdict: DISPROVED.}
Conjecture~00000007675 is \textbf{false}. Specifically:
\begin{enumerate}
  \item \textbf{Identically vanishing valuation.} For every prime $p$, every integer
    $r \ge 1$, and all integers $0 \le k \le n$, we have
    $\qbinom{n}{k}{p^r} \equiv 1 \pmod{p^r}$ and therefore
    $v_p\!\left(\qbinom{n}{k}{p^r}\right) = 0$ (Theorem~\ref{thm:vanishing}).
  \item \textbf{Clause~1 (carry-difference sum).} Under the standard $1$-indexed
    base-$p$ carry convention ($\carry_j(k,n-k)$ being the carry into $p^j$),
    Kummer's termwise identity implies that every summand in Clause~1 is
    identically $0$, so Clause~1 reduces to the trivial identity $0 = 0$
    (Corollary~\ref{cor:clause1-1idx}). Under $0$-indexed carries ($\carry_j$ out of
    $p^j$ into $p^{j+1}$, summed over $j \ge 1$), Clause~1 already fails at
    $(p,n,k) = (2,2,1)$ (Proposition~\ref{prop:clause1-0idx}).
  \item \textbf{Clause~2 (alternative LCM-expression ``holds only when $p \le n - k$'').}
    Clause~2 is false under every conceivable parsing and quantifier scope:
    \begin{itemize}
      \item \emph{Refutation of necessity (``only when $p \le n - k$''):} For
        \textbf{every odd prime $p \ge 3$}, at the strictly interior pair
        $(n,k) = (2,1)$ we have $n - k = 1 < p$ (so $p \le n - k$ is \emph{false}),
        yet $\lcm(2)/ (1!\,1!) = 2 \in [1,p)$ gives
        $\lfloor \log_p(2/1)\rfloor = 0 = v_p\!\left(\qbinom{2}{1}{p}\right)$, and
        in fact all eight natural parsings $\mathrm{LCM}_1,\dots,\mathrm{LCM}_8$
        evaluate to $0$ (Theorem~\ref{thm:odd-family}). At $p = 2$, the strictly
        interior pair $(n,k) = (3,2)$ likewise satisfies $n - k = 1 < 2 = p$ while
        $\lfloor \log_2(\lcm(3)/(2!\,1!))\rfloor = \lfloor \log_2(3/2)\rfloor = 0
        = v_2\!\left(\qbinom{3}{2}{2}\right)$ (Theorem~\ref{thm:per-prime}).
      \item \emph{Refutation of sufficiency (``when $p \le n - k$'') and biconditional:}
        Conversely, when $p \le n - k$ \emph{does} hold with $0 < k < n$, the
        LCM-expression fails to equal $0$: for instance, at $(p,n,k) = (2,4,1)$
        ($p = 2 \le 3 = n - k$) we have $\lcm(2,3,4)/(1!\,3!) = 12/6 = 2$ and
        $\lfloor \log_2 2 \rfloor = 1 \neq 0$; moreover, for \textbf{every prime
        $p \ge 2$}, at $(n,k) = (2p+6,1)$ we have $p \le 2p+5 = n-k$ yet
        $\lcm(2,\dots,2p+6) \le (2p+5)!/2 < 1!\,(2p+5)!$, forcing
        $\lfloor \log_p(\lcm(2,\dots,2p+6)/(1!\,(2p+5)!))\rfloor \le -1 < 0$
        (Theorem~\ref{thm:per-prime} and Theorem~\ref{thm:all-eight}).
    \end{itemize}
\end{enumerate}

\section{Gaussian Binomial Coefficients and Vanishing Valuation}

\begin{definition}[Gaussian binomial coefficients and $q$-factorials]
\label{def:qbinom}
For $q \in \mathbb{N}$ and $m \in \mathbb{N}$, define the $q$-integer and
$q$-factorial by
\[
  [0]_q := 0, \qquad [m+1]_q := [m]_q + q^m = \sum_{i=0}^{m} q^i,
  \qquad
  [0]_q! := 1, \qquad [m+1]_q! := [m]_q! \cdot [m+1]_q.
\]
Define the Gaussian binomial coefficient $\qbinom{n}{k}{q} \in \mathbb{N}$ for
$n, k \in \mathbb{N}$ via the $q$-Pascal recurrence:
\[
  \qbinom{n}{0}{q} := 1, \qquad
  \qbinom{0}{k+1}{q} := 0, \qquad
  \qbinom{n+1}{k+1}{q} := \qbinom{n}{k+1}{q} + q^{n-k} \qbinom{n}{k}{q}.
\]
\end{definition}

By immediate induction on $n$, $\qbinom{n}{k}{q} = 0$ whenever $n < k$, and
$\qbinom{n}{n}{q} = 1$ for all $n \ge 0$. At $q = 1$, the recurrence reduces to
Pascal's identity, so $\qbinom{n}{k}{1} = \binom{n}{k}$. We first record both
standard closed-form characterizations to confirm that Definition~\ref{def:qbinom}
coincides with the canonical Gaussian binomial coefficient for all parameters.

\begin{proposition}[Canonical $q$-factorial and multiplicative identities]
\label{prop:qbinom-canon}
Let $0 \le k \le n$ be integers.
\begin{enumerate}
  \item For every $q \in \mathbb{N}$,
    $[k]_q! \, [n-k]_q! \, \qbinom{n}{k}{q} = [n]_q!$, where $[m]_q! > 0$ whenever
    $q \ge 1$.
  \item For every integer $q \ge 2$, both sides below are positive and satisfy
    \[
      \left(\prod_{i=1}^{k} (q^i - 1)\right) \qbinom{n}{k}{q}
      = \prod_{i=1}^{k} (q^{n-k+i} - 1).
    \]
\end{enumerate}
\end{proposition}

\begin{proof}
First note that $[a]_q + q^a [b]_q = [a+b]_q$ for all $a,b \in \mathbb{N}$ by
induction on $b$: the base case $b=0$ is immediate, and for $b+1$,
$[a]_q + q^a([b]_q + q^b) = ([a]_q + q^a[b]_q) + q^{a+b} = [a+b]_q + q^{a+b}
= [a+b+1]_q$.

We prove (1) by induction on $n$. For $k = 0$ or $k = n$, we have
$\qbinom{n}{k}{q} = 1$ and $[0]_q! = 1$, so equality is immediate. For an inductive
step at $(n+1, k+1)$ with $0 \le k < n$, write $n - k = m + 1 \ge 1$, where
$m = n - (k+1)$. Using Definition~\ref{def:qbinom} and the induction hypotheses at
$(n, k+1)$ and $(n, k)$,
\begin{align*}
  [k+1]_q!\,[n-k]_q!\,\qbinom{n+1}{k+1}{q}
  &= [m+1]_q \bigl([k+1]_q!\,[m]_q!\,\qbinom{n}{k+1}{q}\bigr) \\
  &\quad + q^{n-k}[k+1]_q \bigl([k]_q!\,[n-k]_q!\,\qbinom{n}{k}{q}\bigr) \\
  &= \bigl([n-k]_q + q^{n-k}[k+1]_q\bigr)[n]_q!
   = [n+1]_q\,[n]_q! = [n+1]_q!.
\end{align*}

For (2), let $q \ge 2$. Since $(q-1)[m]_q = q^m - 1$ for all $m \ge 0$, induction
on $m$ gives $(q-1)^m [m]_q! = \prod_{i=1}^{m} (q^i - 1)$. Multiplying the identity
in (1) by $(q-1)^n = (q-1)^k (q-1)^{n-k}$ yields
\[
  \left(\prod_{i=1}^{n-k} (q^i - 1)\right)
  \left(\prod_{i=1}^{k} (q^i - 1)\right) \qbinom{n}{k}{q}
  = \prod_{i=1}^{n} (q^i - 1)
  = \left(\prod_{i=1}^{n-k} (q^i - 1)\right)
    \left(\prod_{i=1}^{k} (q^{n-k+i} - 1)\right).
\]
Since $q \ge 2$, the factor $\prod_{i=1}^{n-k}(q^i - 1)$ is strictly positive and
cancels.
\end{proof}

\begin{definition}[$p$-adic valuation]
For an integer $p \ge 2$ and a positive integer $x \ge 1$, the $p$-adic valuation
$v_p(x) \in \mathbb{N}$ is the unique non-negative integer $v$ such that
$p^v \mid x$ and $p^{v+1} \nmid x$.
\end{definition}

\begin{theorem}[Congruence modulo $q$ and identically zero prime-power valuation]
\label{thm:vanishing}
\leavevmode
\begin{enumerate}
  \item For every $q \in \mathbb{N}$ and all integers $0 \le k \le n$, there exists
    $m \in \mathbb{N}$ such that $\qbinom{n}{k}{q} = 1 + q m$.
  \item For every integer $p \ge 2$, every multiple $q$ of $p$, and all
    $0 \le k \le n$, we have $p \nmid \qbinom{n}{k}{q}$ and hence
    $v_p\!\left(\qbinom{n}{k}{q}\right) = 0$.
  \item In particular, for every prime $p$, every integer $r \ge 1$, and all
    $0 \le k \le n$,
    \[
      v_p\!\left(\qbinom{n}{k}{p^r}\right) = 0
      \qquad\text{and}\qquad
      v_p\!\left(\qbinom{n}{k}{p}\right) = 0.
    \]
\end{enumerate}
\end{theorem}

\begin{proof}
We prove (1) by induction on $n$. If $k = 0$ or $k = n$, then
$\qbinom{n}{k}{q} = 1 = 1 + q \cdot 0$. For $(n+1, k+1)$ with $0 \le k < n$, write
$n - k = d + 1$ with $d \ge 0$. By the induction hypothesis on $(n, k+1)$, there
exists $m_1 \in \mathbb{N}$ such that $\qbinom{n}{k+1}{q} = 1 + q m_1$. Then
\[
  \qbinom{n+1}{k+1}{q}
  = (1 + q m_1) + q^{d+1}\qbinom{n}{k}{q}
  = 1 + q\left(m_1 + q^d \qbinom{n}{k}{q}\right),
\]
completing the induction.

For (2), write $q = p t$ with $t \in \mathbb{N}$. By (1),
$\qbinom{n}{k}{q} = 1 + p(t m) \ge 1$. Since $p \ge 2$, no multiple of $p$ can equal
$1 + p(t m)$ (as $p s \le p(t m) < 1 + p(t m) < p(t m + 1) \le p s'$ whenever
$s \le t m < t m + 1 \le s'$). Thus $p^0 = 1 \mid \qbinom{n}{k}{q}$ and
$p^1 \nmid \qbinom{n}{k}{q}$, so $v_p\!\left(\qbinom{n}{k}{q}\right) = 0$.
Part~(3) follows immediately from $p \mid p^r$ for all $r \ge 1$ and the case $r = 1$.
\end{proof}

\begin{remark}[Origin of the mathematical conflation]
\label{rem:knuth-wilf}
In the literature on $q$-binomial divisibility (Fray, 1967; Knuth and Wilf, 1989),
one studies $v_p\!\left(\qbinom{n}{k}{q}\right)$ when $\gcd(q,p) = 1$, where the
cyclotomic factorization $\qbinom{n}{k}{q} = \prod_{d \ge 2}
\Phi_d(q)^{\lfloor n/d\rfloor - \lfloor k/d\rfloor - \lfloor (n-k)/d\rfloor}$ yields
a carry-counting formula governed by the multiplicative order $\operatorname{ord}_p(q)$.
When instead $q = p$ (or $q = p^r$ with $r \ge 1$), $\gcd(q,p) = p \neq 1$; every
cyclotomic polynomial $\Phi_d(X)$ with $d \ge 2$ has constant term $\Phi_d(0) = 1$,
so $\Phi_d(p^r) \equiv 1 \pmod{p}$ and the $p$-adic valuation collapses identically
to $0$.
\end{remark}

\section{Analysis of Clause 1: Base-\texorpdfstring{$p$}{p} Carries and Kummer's Identity}

Let $p \ge 2$ and $a, b \in \mathbb{N}$. For $i \ge 0$, write
$d_i(x) := \lfloor x / p^i \rfloor \bmod p \in \{0,\dots,p-1\}$ for the $i$-th
base-$p$ digit of $x$. The schoolbook ripple-carry sequence is defined by
\[
  c_0(a,b) := 0, \qquad
  c_{j+1}(a,b) := \left\lfloor \frac{d_j(a) + d_j(b) + c_j(a,b)}{p} \right\rfloor
  \quad (j \ge 0),
\]
so for $j \ge 1$, $c_j(a,b)$ is the $j$-th carry (out of the $p^{j-1}$ digit into
the $p^j$ place). Define also the modulo-$p^j$ residue indicator
\[
  r_j(a,b) :=
  \begin{cases}
    1 & \text{if } (a \bmod p^j) + (b \bmod p^j) \ge p^j, \\
    0 & \text{otherwise.}
  \end{cases}
\]

\begin{lemma}[Equivalence of ripple and residue carries]
\label{lem:carry-equiv}
For every integer $p \ge 2$, all $a, b \in \mathbb{N}$, and all $j \ge 0$,
$c_j(a,b) = r_j(a,b) \in \{0,1\}$.
\end{lemma}

\begin{proof}
For $j = 0$, $c_0(a,b) = 0 = r_0(a,b)$ since $(a \bmod 1) + (b \bmod 1) = 0 < 1$.
Suppose $c_j(a,b) = r_j(a,b)$ holds at $j$. By Euclidean division, any $x \in \mathbb{N}$
satisfies $x \bmod p^{j+1} = p^j d_j(x) + (x \bmod p^j)$ because the right-hand side
is congruent to $x$ modulo $p^{j+1}$ and strictly less than
$p^j(d_j(x) + 1) \le p^{j+1}$. Adding this decomposition for $a$ and $b$ and writing
$(a \bmod p^j) + (b \bmod p^j) = p^j r_j(a,b) + s_j$ with $0 \le s_j < p^j$, we obtain
\[
  (a \bmod p^{j+1}) + (b \bmod p^{j+1})
  = p^j \bigl(d_j(a) + d_j(b) + r_j(a,b)\bigr) + s_j.
\]
Since $0 \le s_j < p^j$, the left-hand side is at least $p^{j+1} = p^j \cdot p$ if
and only if $d_j(a) + d_j(b) + r_j(a,b) \ge p$. As
$d_j(a) + d_j(b) + r_j(a,b) \le (p-1) + (p-1) + 1 = 2p - 1 < 2p$, the quotient
$c_{j+1}(a,b) = \lfloor(d_j(a) + d_j(b) + c_j(a,b))/p\rfloor$ equals $1$ in that case
and $0$ otherwise, matching $r_{j+1}(a,b)$.
\end{proof}

\begin{theorem}[Kummer's termwise floor--carry identity]
\label{thm:kummer-term}
For every integer $p \ge 2$, all $a, b \in \mathbb{N}$, and every $j \ge 0$,
\[
  \left\lfloor \frac{a+b}{p^j} \right\rfloor
  = \left\lfloor \frac{a}{p^j} \right\rfloor
  + \left\lfloor \frac{b}{p^j} \right\rfloor
  + c_j(a,b).
\]
\end{theorem}

\begin{proof}
Let $m = p^j \ge 1$, and write $a = m q_a + u_a$ and $b = m q_b + u_b$ with
$0 \le u_a, u_b < m$. By Lemma~\ref{lem:carry-equiv}, $c_j(a,b) = r_j(a,b)$.
If $u_a + u_b < m$, then $a + b = m(q_a + q_b) + (u_a + u_b)$ has remainder
$u_a + u_b \in [0, m)$, so $\lfloor (a+b)/m \rfloor = q_a + q_b + 0$.
If $u_a + u_b \ge m$, then $0 \le u_a + u_b - m < m$ and
$a + b = m(q_a + q_b + 1) + (u_a + u_b - m)$, so
$\lfloor (a+b)/m \rfloor = q_a + q_b + 1$.
\end{proof}

\begin{corollary}[Clause~1 under $1$-indexed carries]
\label{cor:clause1-1idx}
Let $p$ be a prime and $0 \le k \le n$. Setting $\carry_j(k,n-k) := c_j(k,n-k)$ for
$j \ge 1$, every term of the series in Clause~1 satisfies
\[
  \left\lfloor \frac{n}{p^j} \right\rfloor
  - \left\lfloor \frac{k}{p^j} \right\rfloor
  - \left\lfloor \frac{n-k}{p^j} \right\rfloor
  - \carry_j(k,n-k) = 0,
\]
with all four components individually equal to $0$ whenever $p^j > n$. Hence the
series is finitely supported and equals $0 = v_p\!\left(\qbinom{n}{k}{p}\right)$
tautologically.
\end{corollary}

\begin{proof}
Apply Theorem~\ref{thm:kummer-term} with $a = k$ and $b = n - k$ ($a + b = n$) and
compare with Theorem~\ref{thm:vanishing}(3).
\end{proof}

\begin{proposition}[Failure of Clause~1 under $0$-indexed carries]
\label{prop:clause1-0idx}
If instead $\carry_j(k,n-k)$ is interpreted $0$-index-wise as the carry out of the
$p^j$ place into $p^{j+1}$, i.e., $\carry_j^{(0)}(k,n-k) := c_{j+1}(k,n-k)$, then at
$(p,n,k) = (2,2,1)$ the Clause~1 series equals $1 \neq 0
= v_2\!\left(\qbinom{2}{1}{2}\right)$.
\end{proposition}

\begin{proof}
At $(p,n,k) = (2,2,1)$, $\qbinom{2}{1}{2} = 1 + 2 = 3$, so
$v_2\!\left(\qbinom{2}{1}{2}\right) = 0$. Meanwhile, for $j = 1$,
$\lfloor 2/2 \rfloor - \lfloor 1/2 \rfloor - \lfloor 1/2 \rfloor - c_2(1,1)
= 1 - 0 - 0 - 0 = 1$, while for every $j \ge 2$ we have $2^j \ge 4 > 2$ and
$2^{j+1} \ge 8 > 2$, so every component for $j \ge 2$ vanishes. Thus the sum equals
$1 \neq 0$.
\end{proof}

\section{Disproof of Clause 2 and of the Full Conjecture}

We now examine Clause~2, which asserts that $v_p\!\left(\qbinom{n}{k}{p}\right)$
equals the alternative LCM-expression \emph{only when} $p \le n - k$.

\subsection{Exact Specifications and All Eight Candidate Parsings}

For integers $1 \le a \le b$, write $L(a,b) := \lcm(a, a+1, \dots, b) \ge 1$ (with
the standard empty convention $L(a,b) = 1$ if $a > b$; all counterexamples below satisfy
$0 < k < n$, so $k + 1 \le n$ and the interval $\{k+1,\dots,n\}$ is non-empty).
For an integer $p \ge 2$ and positive integers $u, v \ge 1$, the exact floor
base-$p$ logarithm $e = \lfloor \log_p(u/v) \rfloor \in \mathbb{Z}$ is uniquely
determined by $p^e \le u/v < p^{e+1}$, i.e.,
\[
  \begin{cases}
    v p^m \le u < v p^{m+1} & \text{if } e = m \ge 0,\\[2pt]
    v \le u p^{m+1} \text{ and } u p^m < v & \text{if } e = -(m+1) < 0.
  \end{cases}
\]
Likewise, for any integer $L \ge 1$ and $d \ge 1$, letting
$q = \lfloor \lfloor \log_p L \rfloor / d \rfloor$, Euclidean division
$d q \le \lfloor \log_p L \rfloor < d(q+1)$ yields $p^{d q} \le L < p^{d(q+1)}$, so
$\lfloor (\log_p L) / d \rfloor = \lfloor \lfloor \log_p L \rfloor / d \rfloor$ over
$\mathbb{R}$.

To ensure that no syntactic or semantic reading of the phrase
``$\lfloor \log_p \lcm(k+1,\dots,n)/(k!(n-k)!)\rfloor$ type expression'' escapes our
disproof, Table~\ref{tab:parsings} lists all eight natural candidate expressions
$\mathrm{LCM}_1,\dots,\mathrm{LCM}_8$.

\begin{table}[ht]
\centering
\small
\renewcommand{\arraystretch}{1.35}
\begin{tabular}{@{}ll>{\raggedright\arraybackslash}p{6.2cm}@{}}
\hline
\textbf{Label} & \textbf{Mathematical Formula} & \textbf{Interpretation} \\
\hline
$\mathrm{LCM}_1$ & $\displaystyle\left\lfloor \log_p \frac{\lcm(k+1,\dots,n)}{k!\,(n-k)!} \right\rfloor$
  & Primary literal parsing ($\mathbb{Q}_{>0}$ quotient inside $\log_p$) \\
$\mathrm{LCM}_2$ & $\displaystyle\left\lfloor \frac{\log_p \lcm(k+1,\dots,n)}{k!\,(n-k)!} \right\rfloor$
  & Outer-quotient parsing ($\log_p \lcm$ divided by $k!(n-k)!$) \\
$\mathrm{LCM}_3$ & $\displaystyle\left\lfloor \log_p \frac{\lcm(k+1,\dots,n)}{(n-k)!} \right\rfloor$
  & Reduced denominator $(n-k)!$ inside $\log_p$ \\
$\mathrm{LCM}_4$ & $v_p(\lcm(k+1,\dots,n)) - v_p(k!\,(n-k)!)$
  & Exact $p$-adic valuation of $\lcm(k+1,\dots,n)/(k!(n-k)!)$ \\
$\mathrm{LCM}_5$ & $v_p(\lcm(k+1,\dots,n)) - v_p((n-k)!)$
  & Exact $p$-adic valuation of $\lcm(k+1,\dots,n)/(n-k)!$ \\
$\mathrm{LCM}_6$ & $\lfloor \log_p \lcm(k+1,\dots,n)\rfloor - v_p(k!\,(n-k)!)$
  & Hybrid $\lfloor \log_p \lcm \rfloor$ minus full factorial valuation \\
$\mathrm{LCM}_7$ & $\lfloor \log_p \lcm(k+1,\dots,n)\rfloor - v_p((n-k)!)$
  & Hybrid $\lfloor \log_p \lcm \rfloor$ minus $(n-k)!$ valuation \\
$\mathrm{LCM}_8$ & $v_p\!\left(\binom{n}{k}\right)$
  & Classical binomial $p$-adic valuation \\
\hline
\end{tabular}
\caption{All eight candidate parsings of the Clause~2 LCM-expression.}
\label{tab:parsings}
\end{table}

\subsection{Universal Infinite-Family and Per-Prime Refutations}

\begin{theorem}[Infinite-family refutation of ``only when $p \le n - k$'' across all eight parsings]
\label{thm:odd-family}
Let $p \ge 3$ be any odd prime, and consider the strictly interior pair
$(n,k) = (2,1)$ ($0 < k < n$). Then:
\begin{enumerate}
  \item $n - k = 1 < 3 \le p$, so $p \le n - k$ is \textbf{false};
  \item $v_p\!\left(\qbinom{2}{1}{p}\right) = v_p(1 + p) = 0$; and
  \item $\mathrm{LCM}_m(p,2,1) = 0 = v_p\!\left(\qbinom{2}{1}{p}\right)$
    simultaneously for \textbf{all eight} parsings $m \in \{1,\dots,8\}$.
\end{enumerate}
Consequently, the literal assertion that the LCM-expression holds \emph{only when}
$p \le n - k$ fails across an infinite family of primes under every parsing
$\mathrm{LCM}_1,\dots,\mathrm{LCM}_8$.
\end{theorem}

\begin{proof}
At $(n,k) = (2,1)$, the interval $\{k+1,\dots,n\} = \{2\}$ is a singleton, so
$\lcm(2,\dots,2) = 2$, while $k! = 1! = 1$, $(n-k)! = 1! = 1$, $k!(n-k)! = 1$, and
$\binom{2}{1} = 2$. Since $p \ge 3$, we have $p^0 = 1 \le 2 < 3 \le p^1$, giving
$\lfloor \log_p(2/1)\rfloor = 0$ and $\lfloor \log_p 2\rfloor = 0$. Likewise $p > 2$
implies $p \nmid 2$ and $p \nmid 1$, so $v_p(2) = 0$ and $v_p(1) = 0$. Direct
substitution into Table~\ref{tab:parsings} gives $\mathrm{LCM}_m(p,2,1) = 0$ for all
$1 \le m \le 8$, matching $v_p\!\left(\qbinom{2}{1}{p}\right) = 0$ from
Theorem~\ref{thm:vanishing}(3), whereas $p \le 2 - 1 = 1$ is false.
\end{proof}

Next we show that even if $p$ is treated as a \emph{fixed} prime parameter (including
$p = 2$), the primary literal expression $\mathrm{LCM}_1$ fails \emph{both} directions
of $p \le n - k$ on strictly interior parameters $0 < k < n$ at that same prime~$p$.

\begin{theorem}[Per-prime two-sided refutation for every prime $p$]
\label{thm:per-prime}
Let $p$ be \textbf{any} prime ($p = 2$ or $p \ge 3$).
\begin{enumerate}
  \item \emph{(Failure of ``only when $p \le n - k$'')} There exist integers
    $0 < k < n$ with $p > n - k$ such that
    $\mathrm{LCM}_1(p,n,k) = \mathrm{LCM}_2(p,n,k) = 0
    = v_p\!\left(\qbinom{n}{k}{p}\right)$.
  \item \emph{(Failure of ``when $p \le n - k$'')} At $(n,k) = (2p+6, 1)$, we have
    $0 < k < n$ and $p \le 2p + 5 = n - k$, yet
    \[
      \mathrm{LCM}_1(p, 2p+6, 1) \le -1 < 0 = v_p\!\left(\qbinom{2p+6}{1}{p}\right).
    \]
\end{enumerate}
\end{theorem}

\begin{proof}
For (1), if $p \ge 3$, take $(n,k) = (2,1)$ and apply Theorem~\ref{thm:odd-family}.
If $p = 2$, take $(n,k) = (3,2)$: then $0 < 2 < 3$ and $n - k = 1 < 2 = p$, while
$\lcm(3) = 3$ and $2!\,1! = 2$. Since $2 \cdot 2^0 = 2 \le 3 < 4 = 2 \cdot 2^1$, we
have $\mathrm{LCM}_1(2,3,2) = \lfloor \log_2(3/2)\rfloor = 0$, and since
$\lfloor \log_2 3 \rfloor = 1$, we also have
$\mathrm{LCM}_2(2,3,2) = \lfloor 1/2 \rfloor = 0$, both matching
$v_2\!\left(\qbinom{3}{2}{2}\right) = 0$.

For (2), let $p \ge 2$ be arbitrary and set $n = 2p + 6$, $k = 1$, so $0 < 1 < n$
and $n - k = 2p + 5 \ge p$. Define the positive integer
\[
  M := \prod_{j=3}^{2p+5} j = \frac{(2p+5)!}{2}.
\]
Because $p \ge 2$, we have the strict chain of inequalities
$3 \le 4 < 5 \le p + 3 \le 2p + 5$. Hence $4$ and $p + 3$ are two \emph{distinct}
factors occurring in the product $M$, which implies $4(p+3) \mid M$. Every integer
$m \in \{2, 3, \dots, 2p+6\}$ divides $M$:
\begin{itemize}
  \item if $m = 2$, then $2 \mid 4 \mid 4(p+3) \mid M$;
  \item if $3 \le m \le 2p+5$, then $m$ is one of the factors defining $M$;
  \item if $m = 2p + 6 = 2(p+3)$, then $2(p+3) \mid 4(p+3) \mid M$.
\end{itemize}
By the universal property of the least common multiple,
$\lcm(2,3,\dots,2p+6) \mid M$. Since $M > 0$, this gives
\[
  \lcm(2,3,\dots,2p+6) \le M < 2M = (2p+5)! = 1!\,(n-k)!.
\]
Therefore $\frac{\lcm(2,\dots,2p+6)}{1!\,(2p+5)!} < 1 = p^0$, so
$\mathrm{LCM}_1(p,2p+6,1) \le -1 < 0 = v_p\!\left(\qbinom{2p+6}{1}{p}\right)$.
\end{proof}

\begin{theorem}[Complete refutation of all eight parsings in both directions]
\label{thm:all-eight}
For every $m \in \{1,\dots,8\}$, on strictly interior parameters $0 < k < n$:
\begin{enumerate}
  \item \emph{Necessity fails:} There exist a prime $p$ and $0 < k < n$ with
    $\mathrm{LCM}_m(p,n,k) = v_p\!\left(\qbinom{n}{k}{p}\right)$ and $p > n - k$.
  \item \emph{Sufficiency fails:} There exist a prime $p$ and $0 < k < n$ with
    $p \le n - k$ and $\mathrm{LCM}_m(p,n,k) \neq v_p\!\left(\qbinom{n}{k}{p}\right)$.
  \item \emph{Biconditional fails:} In particular,
    $\mathrm{LCM}_m(p,n,k) = v_p\!\left(\qbinom{n}{k}{p}\right) \iff p \le n - k$
    is false in both directions for every $m \in \{1,\dots,8\}$.
\end{enumerate}
\end{theorem}

\begin{proof}
Part~(1) is established simultaneously for all $1 \le m \le 8$ at $(p,n,k) = (3,2,1)$
by Theorem~\ref{thm:odd-family}. Part~(2) is established by direct arithmetic on the
two interior triples $(p,n,k) = (2,3,1)$ ($n - k = 2 \ge 2 = p$, $\lcm(2,3) = 6$,
$1!\,2! = 2$) and $(p,n,k) = (2,4,1)$ ($n - k = 3 \ge 2 = p$, $\lcm(2,3,4) = 12$,
$1!\,3! = 6$, $\binom{4}{1} = 4$), summarized in Table~\ref{tab:witnesses}:
at $(2,3,1)$, $\mathrm{LCM}_2(2,3,1) = \lfloor \log_2 6 \rfloor / 2 = 2/2 = 1 \neq 0$;
at $(2,4,1)$, $\mathrm{LCM}_1 = \mathrm{LCM}_3 = \lfloor \log_2(12/6)\rfloor = 1$,
$\mathrm{LCM}_4 = \mathrm{LCM}_5 = v_2(12) - v_2(6) = 2 - 1 = 1$,
$\mathrm{LCM}_6 = \mathrm{LCM}_7 = \lfloor \log_2 12\rfloor - v_2(6) = 3 - 1 = 2$,
and $\mathrm{LCM}_8 = v_2(4) = 2$, all non-zero whereas
$v_2\!\left(\qbinom{n}{k}{2}\right) = 0$. Part~(3) follows immediately from (1)
and~(2).
\end{proof}

\begin{table}[ht]
\centering
\footnotesize
\setlength{\tabcolsep}{4.5pt}
\renewcommand{\arraystretch}{1.25}
\begin{tabular}{@{}cccccccccccc@{}}
\hline
$(p,n,k)$ & $n-k$ & $p \le n-k$ & $v_p\!\left(\qbinom{n}{k}{p}\right)$
  & $\mathrm{LCM}_1$ & $\mathrm{LCM}_2$ & $\mathrm{LCM}_3$ & $\mathrm{LCM}_4$
  & $\mathrm{LCM}_5$ & $\mathrm{LCM}_6$ & $\mathrm{LCM}_7$ & $\mathrm{LCM}_8$ \\
\hline
$(3,2,1)$ & $1$ & False & $0$ & $0$ & $0$ & $0$ & $0$ & $0$ & $0$ & $0$ & $0$ \\
$(2,3,2)$ & $1$ & False & $0$ & $0$ & $0$ & $1$ & $-1$ & $0$ & $0$ & $1$ & $0$ \\
$(2,3,1)$ & $2$ & True  & $0$ & $1$ & $1$ & $1$ & $0$ & $0$ & $1$ & $1$ & $0$ \\
$(2,4,1)$ & $3$ & True  & $0$ & $1$ & $0$ & $1$ & $1$ & $1$ & $2$ & $2$ & $2$ \\
$(3,6,1)$ & $5$ & True  & $0$ & $-1$ & $0$ & $-1$ & $0$ & $0$ & $2$ & $2$ & $1$ \\
\hline
\end{tabular}
\caption{Exact values at representative strictly interior counterexamples ($0 < k < n$).}
\label{tab:witnesses}
\end{table}

\begin{corollary}[Disproof of Conjecture 00000007675]
\label{cor:main-disproof}
Conjecture~00000007675 is false.
\end{corollary}

\begin{remark}[What if $q = 1$ had been substituted instead of $q = p$?]
Even if one disregarded the explicit definition ``substituting $q = p$'' and replaced
$\qbinom{n}{k}{p}$ by the classical binomial coefficient
$\qbinom{n}{k}{1} = \binom{n}{k}$, \emph{both} clauses of Conjecture~00000007675
would still fail: Clause~1 evaluates to $0$ (by Corollary~\ref{cor:clause1-1idx})
whereas $v_2\!\binom{2}{1} = 1 \neq 0$, and Clause~2 still fails at $(p,n,k) = (3,2,1)$
because $v_3\!\binom{2}{1} = v_3(2) = 0 = \mathrm{LCM}_m(3,2,1)$ despite
$p = 3 > 1 = n - k$.
\end{remark}

\section{Formal Verification in Core Lean 4.19.0 and Computational Audit}

All definitions and proofs in Sections~2--4 are formalized from first principles in
\texttt{lean/Main.lean} in Core Lean~4.19.0 (\texttt{leanprover/lean4:v4.19.0}),
compiled with strict warning-as-error checking and zero external library
dependencies. Every computable primitive (\texttt{padicVal}, \texttt{qBinom},
\texttt{carryRipple}, \texttt{lcmInterval}, and \texttt{floorLogQ}) is paired with
proved existence and uniqueness theorems guaranteeing exact semantic fidelity.
Table~\ref{tab:lean-map} maps our mathematical results to their verified
declarations in \texttt{Main.lean}.

\begin{table}[ht]
\centering
\footnotesize
\setlength{\tabcolsep}{5pt}
\renewcommand{\arraystretch}{1.25}
\begin{tabular}{@{}ll@{}}
\hline
\textbf{Paper Result} & \textbf{Core Lean 4.19.0 Declaration (\texttt{Main.lean})} \\
\hline
Spec \& uniqueness of $v_p(x)$ & \texttt{padicVal\_spec}, \texttt{isPadicVal\_unique} \\
Proposition~\ref{prop:qbinom-canon}(1)--(2) & \texttt{qBinom\_mul\_qFact}, \texttt{qFallProd\_mul\_qBinom} \\
Theorem~\ref{thm:vanishing}(1)--(3) & \texttt{qBinom\_eq\_one\_add\_mul}, \texttt{padicVal\_qBinom\_prime\_pow\_eq\_zero} \\
Lemma~\ref{lem:carry-equiv} & \texttt{carryRipple\_eq\_carryResidue} \\
Theorem~\ref{thm:kummer-term} \& Corollary~\ref{cor:clause1-1idx} & \texttt{kummer\_div\_add\_carry}, \texttt{clause1\_holds\_under\_one\_indexed\_carry} \\
Proposition~\ref{prop:clause1-0idx} & \texttt{clause1\_fails\_under\_zero\_indexed\_carry} \\
Universal spec of $\lcm(a,\dots,b)$ & \texttt{lcmInterval\_spec}, \texttt{isIntervalLCM\_unique} \\
Spec \& uniqueness of $\lfloor \log_p(u/v)\rfloor$ & \texttt{floorLogQ\_spec}, \texttt{isFloorLogQ\_unique}, \texttt{floorLogNat\_div\_exact\_spec} \\
Theorem~\ref{thm:odd-family} & \texttt{odd\_prime\_refutes\_only\_when\_all\_eight} \\
Theorem~\ref{thm:per-prime} & \texttt{every\_prime\_refutes\_both\_directions\_lcmExpr1} \\
Theorem~\ref{thm:all-eight} & \texttt{all\_eight\_parsings\_refuted\_both\_directions} \\
Corollary~\ref{cor:main-disproof} & \texttt{conjecture\_00000007675\_disproved} \\
Non-vacuity \& concrete checks & \texttt{concrete\_witness\_checks} \\
\hline
\end{tabular}
\caption{Correspondence between mathematical theorems and verified Lean~4 declarations.}
\label{tab:lean-map}
\end{table}

\noindent
Running \texttt{\#print axioms} on every theorem listed in Table~\ref{tab:lean-map}
confirms that they depend solely on Lean's standard foundational axioms
(\texttt{propext}, \texttt{Classical.choice}, \texttt{Quot.sound}). In addition, the
standalone script \texttt{verify.py} independently cross-checks all three Gaussian
binomial definitions, prime-power valuations, carry conventions, and all eight LCM
parsings across finite parameter grids and infinite prime families.

\begin{thebibliography}{9}

\bibitem{Fray1967}
R.~D.~Fray,
\emph{Congruence properties of ordinary and $q$-binomial coefficients},
Duke Mathematical Journal \textbf{34} (1967), 467--480.

\bibitem{KnuthWilf1989}
D.~E.~Knuth and H.~S.~Wilf,
\emph{The power of a prime that divides a generalized binomial coefficient},
Journal f\"ur die reine und angewandte Mathematik \textbf{396} (1989), 212--219.

\bibitem{Kummer1852}
E.~E.~Kummer,
\emph{\"Uber die Erg\"anzungss\"atze zu den allgemeinen Reciprocit\"atsgesetzen},
Journal f\"ur die reine und angewandte Mathematik \textbf{44} (1852), 93--146.

\end{thebibliography}

\end{document}
